\documentclass[aspectratio=169]{beamer}
\usepackage[utf8]{inputenc}
\usepackage[spanish]{babel}
\usepackage{listings}
\usepackage{tabularx}
\usepackage{booktabs}
\usepackage{colortbl}

% Paquete de estilos personalizados Icesi / IasLab
\usepackage{custombeamer}

% Definición de colores adicionales para tablas y código
\definecolor{cbBorder}{HTML}{D1D5DB}
\definecolor{cbRowLight}{HTML}{F3F4F6}

% Configuración de Listings para YAML y Bash
\lstdefinestyle{yamllst}{
    backgroundcolor=\color{cbLight},
    basicstyle=\ttfamily\scriptsize,
    keywordstyle=\color{cbBlue}\bfseries,
    commentstyle=\color{cbMuted}\itshape,
    stringstyle=\color{cbOrange},
    breaklines=true,
    frame=single,
    rulecolor=\color{cbBorder},
    showstringspaces=false,
    tabsize=2
}

\lstdefinestyle{yamltiny}{
    backgroundcolor=\color{cbLight},
    basicstyle=\ttfamily\tiny,
    keywordstyle=\color{cbBlue}\bfseries,
    commentstyle=\color{cbMuted}\itshape,
    stringstyle=\color{cbOrange},
    breaklines=true,
    frame=single,
    rulecolor=\color{cbBorder},
    showstringspaces=false,
    tabsize=2
}

\begin{document}

% ==============================================================================
% PORTADA
% ==============================================================================
\titleSlideDarkMode
  {GitHub Actions \& CI/CD}
  {Workflows, Jobs, Steps, Runners Autohosteados y Despliegue}
  {Prof. Alejandro Muñoz Bravo}
  {Ingeniería de Software V $\bullet$ Universidad ICESI}

% ==============================================================================
% AGENDA
% ==============================================================================
\begin{frame}{Contenido de la Sesión}
  \vspace{0.2cm}
  \begin{enumerate}[itemsep=10pt]
    \item \textbf{\color{cbBlue}01.} Fundamentos y Arquitectura de GitHub Actions
    \item \textbf{\color{cbBlue}02.} Jerarquía de Componentes: Triggers, Jobs, Stages y Steps
    \item \textbf{\color{cbBlue}03.} Comandos Shell (\texttt{run:}) vs. Acciones Reutilizables (\texttt{uses:})
    \item \textbf{\color{cbBlue}04.} Comparativa Exhaustiva: \textbf{GitHub-Hosted vs. Self-Hosted Runners}
    \item \textbf{\color{cbBlue}05.} Guía Paso a Paso: Construcción del Pipeline CI/CD del Taller
  \end{enumerate}
\end{frame}

% ==============================================================================
% SECCIÓN 1: FUNDAMENTOS
% ==============================================================================
\sectionSlideSolid{01}{¿Qué es GitHub Actions?}{Motor de Automatización y CI/CD Integrado}

\contentSlideBulletCard
  {¿Qué es GitHub Actions?}
  {
    \begin{itemize}
      \item \textbf{Plataforma CI/CD Nativa:} Integrada directamente en el repositorio sin servidores externos.
      \item \textbf{Basada en Eventos:} Reacciona a sucesos del ciclo Git (\texttt{push}, \texttt{pull\_request}, etc.).
      \item \textbf{Runners Aislados:} VMs efímeras provistas por GitHub (\texttt{ubuntu-latest}, \texttt{windows-latest}, \texttt{macos-latest}).
      \item \textbf{Marketplace:} Miles de acciones oficiales y comunitarias listas para reutilizar.
    \end{itemize}
  }
  {Principio Clave}
  {
    \textbf{Shift-Left Testing:}\\[0.15cm]
    Detectar y corregir defectos en etapas tempranas mediante builds y tests automáticos en cada \textit{commit}.
  }

\contentSlideTwoCols
  {Arquitectura y Flujo de Ejecución}
  {Entidades Principales}
  {
    \begin{itemize}
      \item \textbf{Workflow:} Archivo YAML en \texttt{.github/workflows/}.
      \item \textbf{Events (on):} Disparadores del flujo.
      \item \textbf{Jobs:} Tareas en un runner independiente.
      \item \textbf{Steps:} Pasos atómicos secuenciales.
    \end{itemize}
  }
  {Ciclo de Vida}
  {
    \begin{enumerate}
      \item Evento ocurre en Git (\textit{Push} o \textit{PR}).
      \item GitHub asigna una Máquina Virtual (\textit{Runner}).
      \item Se ejecutan los Jobs (paralelo o secuencial).
      \item Se reporta el estado (\textit{Pass / Fail}).
    \end{enumerate}
  }

% ==============================================================================
% SECCIÓN 2: JERARQUÍA DE COMPONENTES
% ==============================================================================
\sectionSlideSolid{02}{Jerarquía de Componentes}{Triggers, Jobs, Stages y Steps}

\begin{frame}[fragile]{1. Triggers y Eventos (\texttt{on:})}
  \vspace{0.1cm}
  El bloque \textbf{\texttt{on:}} define \textit{cuándo} debe dispararse la ejecución del workflow.
  \vspace{0.15cm}
  \begin{lstlisting}[style=yamllst]
name: Continuous Integration & Delivery

on:
  push:
    branches: [ main, develop ]
  pull_request:
    branches: [ main ]
  workflow_dispatch: # Permite disparo manual desde la UI web
  \end{lstlisting}
  \vspace{0.1cm}
  \begin{itemize}
    \item \textbf{Filtros de rama:} Limita la ejecución a ramas estratégicas (\texttt{main}, \texttt{develop}).
    \item \textbf{Filtros por ruta (\texttt{paths}):} Dispara el job solo si cambian carpetas específicas.
  \end{itemize}
\end{frame}

\begin{frame}[fragile]{2. Jobs y Runners (\texttt{jobs: <job\_id>})}
  \vspace{0.1cm}
  Cada \textbf{Job} corre en una máquina virtual limpia. Por defecto, se ejecutan en \textbf{paralelo}.
  \vspace{0.15cm}
  \begin{lstlisting}[style=yamllst]
jobs:
  test-backend:
    runs-on: ubuntu-latest
    steps:
      - uses: actions/checkout@v4
      - name: Run Maven Tests
        run: cd backend && mvn clean test
  \end{lstlisting}
  \vspace{0.1cm}
  \begin{itemize}
    \item \textbf{\texttt{runs-on:}} Especifica el SO del runner (\texttt{ubuntu-latest} es el estándar).
    \item \textbf{Aislamiento total:} Los archivos de un job \textbf{no} existen en otro job a menos que se compartan como artefactos.
  \end{itemize}
\end{frame}

\begin{frame}[fragile]{3. Dependencias entre Jobs (\texttt{needs:})}
  \vspace{0.1cm}
  Para estructurar un pipeline secuencial (\textit{Stages}), se utiliza la directiva \textbf{\texttt{needs:}}.
  \vspace{0.15cm}
  \begin{lstlisting}[style=yamllst]
jobs:
  test-backend:
    runs-on: ubuntu-latest
    ...
  test-frontend:
    runs-on: ubuntu-latest
    ...
  build-docker:
    needs: [test-backend, test-frontend] # Espera a que ambos terminen OK
    runs-on: ubuntu-latest
    ...
  \end{lstlisting}
  \vspace{0.1cm}
  \textbf{Regla de Oro:} Si cualquier job predecesor falla, los jobs dependientes se \textbf{cancelan automáticamente}.
\end{frame}

\begin{frame}[fragile]{4. Steps (\texttt{steps:})}
  \vspace{0.1cm}
  Un \textbf{Step} es una tarea ejecutable dentro del Job. Todos los steps de un job comparten el mismo entorno y disco.
  \vspace{0.15cm}
  \begin{lstlisting}[style=yamllst]
    steps:
      - name: Descargar codigo fuente
        uses: actions/checkout@v4

      - name: Configurar Java 17 Temurin
        uses: actions/setup-java@v4
        with:
          java-version: '17'
          distribution: 'temurin'
          cache: 'maven'

      - name: Ejecutar pruebas unitarias
        run: cd backend && mvn test
  \end{lstlisting}
\end{frame}

% ==============================================================================
% SECCIÓN 3: COMANDOS VS ACCIONES
% ==============================================================================
\sectionSlideSolid{03}{Comandos y Acciones}{run: vs. uses: en la Práctica}

\contentSlideTwoCols
  {run: vs. uses: ¿Cuándo usar cuál?}
  {Comandos Shell (\texttt{run:})}
  {
    \begin{itemize}
      \item Ejecuta comandos directos del CLI del SO (\texttt{bash}, \texttt{sh}).
      \item Ideal para compilar, probar o invocar herramientas locales:
      \begin{itemize}
        \item \texttt{mvn clean test}
        \item \texttt{npm run build}
        \item \texttt{docker compose up -d}
      \end{itemize}
    \end{itemize}
  }
  {Acciones Reutilizables (\texttt{uses:})}
  {
    \begin{itemize}
      \item Módulos empaquetados del Marketplace de GitHub.
      \item Se configuran mediante parámetros en \textbf{\texttt{with:}}.
      \item Acciones esenciales:
      \begin{itemize}
        \item \texttt{actions/checkout@v4}
        \item \texttt{actions/setup-node@v4}
        \item \texttt{docker/build-push-action@v5}
      \end{itemize}
    \end{itemize}
  }

\begin{frame}[fragile]{Compartir Binarios entre Jobs: Upload \& Download Artifacts}
  \vspace{0.1cm}
  Como los runners son máquinas efímeras independientes, usamos artefactos para transferir archivos entre etapas:
  \vspace{0.15cm}
  \begin{lstlisting}[style=yamllst]
# En Job 1: test-backend
- name: Subir artefacto JAR
  uses: actions/upload-artifact@v4
  with:
    name: backend-jar
    path: backend/target/*.jar
    retention-days: 1

# En Job posterior: deploy
- name: Descargar artefacto JAR
  uses: actions/download-artifact@v4
  with:
    name: backend-jar
    path: backend/target/
  \end{lstlisting}
\end{frame}

% ==============================================================================
% SECCIÓN 4: RUNNERS AUTOHOSTEADOS (NUEVA SECCIÓN DETALLADA)
% ==============================================================================
\sectionSlideSolid{04}{GitHub-Hosted vs. Self-Hosted}{Entornos Gestionados vs. Servidores Propios}

\contentSlideBulletCard
  {¿Qué es un Self-Hosted Runner?}
  {
    \begin{itemize}
      \item \textbf{Agente Dedicado:} Software de GitHub (\texttt{actions-runner}) instalado en un servidor propio (ej. clúster IasLab \texttt{grid102}).
      \item \textbf{Acceso a Red Privada:} Puede conectarse directamente a bases de datos internas, LAN o VPN ZeroTier sin abrir puertos a Internet.
      \item \textbf{Hardware a la Medida:} Uso de servidores de alta capacidad (40 vCPUs, 62 GB RAM, GPUs para Deep Learning).
      \item \textbf{Cero Costo de Minutos:} Ejecuciones ilimitadas sin consumir minutos del plan de GitHub.
    \end{itemize}
  }
  {Casos de Uso Típicos}
  {
    \begin{itemize}
      \item Despliegues en servidores de intranet corporativa o universitaria.
      \item Tareas que requieren GPUs o hardware especializado.
      \item Caches masivas en disco local (evita descargas constantes).
    \end{itemize}
  }

\begin{frame}{Matriz Comparativa: GitHub-Hosted vs. Self-Hosted}
  \vspace{-0.1cm}
  \fontsize{7.5}{9.5}\selectfont
  \renewcommand{\arraystretch}{1.12}
  \begin{tabularx}{\textwidth}{|p{2.8cm}|X|X|}
    \hline
    \rowcolor{cbLight}
    \textbf{Dimensión} & \textbf{GitHub-Hosted Runner} & \textbf{Self-Hosted Runner} \\ \hline
    \textbf{Entorno \& Ciclo} & VM efímera (se destruye tras cada job). Disco 100\% limpio. & Host persistente. Workspace y archivos temporales permanecen. \\ \hline
    \textbf{Acceso a Red / LAN} & Solo acceso público en Internet. Requiere VPN/túneles. & Acceso nativo a LAN, VPN (ZeroTier) y servicios internos (\texttt{grid100}, DBs). \\ \hline
    \textbf{Mantenimiento} & GitHub actualiza el SO, parches y herramientas (Node, Java, Docker). & El administrador mantiene SO, Docker daemon, parches y herramientas. \\ \hline
    \textbf{Cómputo / Hardware} & Estándar: 2 vCPU, 7 GB RAM. Hardware predeterminado. & Personalizado (ej. 40 vCPUs, 64 GB RAM, GPUs NVIDIA en \texttt{grid102}). \\ \hline
    \textbf{Costos de Ejecución} & Facturación por minutos consumidos según plan. & Gratuito en minutos de GitHub (costo solo de hardware propio). \\ \hline
  \end{tabularx}
\end{frame}

\begin{frame}[fragile]{Cambios en el YAML del Workflow para Self-Hosted}
  \vspace{-0.1cm}
  Para enrutar un Job a un servidor propio, se modifica la directiva \textbf{\texttt{runs-on:}} utilizando etiquetas (\textit{labels}):
  \vspace{0.05cm}
  \begin{lstlisting}[style=yamltiny]
# Opcion A: Runner estandar de GitHub
jobs:
  deploy-cloud:
    runs-on: ubuntu-latest
    steps:
      - uses: actions/checkout@v4
      - run: mvn clean test

# Opcion B: Runner Autohosteado en Servidor Propio (ej. grid102 IasLab)
jobs:
  deploy-grid100:
    runs-on: [ self-hosted, linux, x64, grid102 ]
    steps:
      - name: Limpiar workspace previo (por persistencia de disco)
        uses: actions/checkout@v4
        with:
          clean: true

      - name: Desplegar directamente en Docker local
        run: docker compose -f docker-compose.prod.yml up -d
  \end{lstlisting}
\end{frame}

\contentSlideTwoCols
  {Seguridad: Consideraciones Críticas en Self-Hosted}
  {Riesgos en Repositorios Públicos}
  {
    \begin{itemize}
      \item \textbf{\color{cbOrange}¡Peligro con Forks y PRs!}\newline Un tercero podría enviar un PR con un script malicioso en el workflow que se ejecute en el servidor interno con acceso a la red de la empresa.
      \item \textbf{Persistencia de Procesos:}\newline Un proceso zombie o un contenedor huérfano puede interferir con jobs futuros.
    \end{itemize}
  }
  {Mejores Prácticas}
  {
    \begin{itemize}
      \item Usar Self-Hosted \textbf{únicamente en repositorios privados} o de confianza interna.
      \item Requerir aprobación manual para ejecuciones provenientes de colaboradores externos.
      \item Utilizar \textbf{Runners Efímeros Contenerizados} (ej. \textit{Actions Runner Controller - ARC} sobre Kubernetes).
    \end{itemize}
  }

% ==============================================================================
% SECCIÓN 5: GUÍA DEL TALLER EVALUATIVO
% ==============================================================================
\sectionSlideSolid{05}{Guía de Construcción del Pipeline}{Estructura de los 4 Jobs del Taller}

\begin{frame}{Estructura General del Pipeline a Construir}
  \centering
  \vspace{0.15cm}
  \begin{tikzpicture}[node distance=0.9cm, auto,
    block/.style={rectangle, draw=cbBlue, fill=cbLight, text=cbDark, thick, rounded corners, minimum height=0.8cm, minimum width=3.8cm, align=center},
    stage/.style={rectangle, draw=cbPurple, fill=cbWhite, text=cbDark, thick, rounded corners, minimum height=0.8cm, minimum width=3.5cm, align=center},
    arrow/.style={thick, ->, >=stealth, color=cbBlue}]

    \node [block, fill=cbBlue, text=cbWhite] (event) {\textbf{Trigger: Push / PR en main}};
    
    \node [stage, below left=0.6cm and -0.6cm of event] (job1) {\textbf{Job 1: test-backend}\\ \tiny Maven Test + JAR Artifact};
    \node [stage, below right=0.6cm and -0.6cm of event] (job2) {\textbf{Job 2: test-frontend}\\ \tiny Node Lint + Vite Build};
    
    \node [stage, below=1.9cm of event, fill=cbOrange!15, draw=cbOrange] (job3) {\textbf{Job 3: build-docker}\\ \tiny Multi-Stage Buildx (API + Web)};
    \node [stage, below=0.6cm of job3, fill=cbGreen!15, draw=cbGreen] (job4) {\textbf{Job 4: smoke-test}\\ \tiny Docker Compose Up \& cURL};

    \draw [arrow] (event) -| (job1);
    \draw [arrow] (event) -| (job2);
    \draw [arrow] (job1) |- (job3);
    \draw [arrow] (job2) |- (job3);
    \draw [arrow] (job3) -- (job4);
  \end{tikzpicture}
\end{frame}

\begin{frame}[fragile]{Paso 1: Configurar Encabezado y Triggers}
  \vspace{0.1cm}
  Crear el archivo \texttt{.github/workflows/ci-cd.yml} y definir el nombre del flujo y los eventos:
  \vspace{0.15cm}
  \begin{lstlisting}[style=yamllst]
name: CI/CD Pipeline - Product Application

on:
  push:
    branches: [ main ]
  pull_request:
    branches: [ main ]

jobs:
  # Aqui definiremos los 4 jobs de forma ordenada
  \end{lstlisting}
  \vspace{0.1cm}
  \textbf{Objetivo:} Asegurar que el pipeline se dispare tanto al hacer \texttt{git push} directo como al abrir un \textit{Pull Request}.
\end{frame}

\begin{frame}[fragile]{Paso 2: Implementar Job 1 (\texttt{test-backend})}
  \vspace{0.05cm}
  Responsable de validar la lógica del microservicio Spring Boot:
  \vspace{0.1cm}
  \begin{lstlisting}[style=yamltiny]
  test-backend:
    name: Backend Unit Tests & Package
    runs-on: ubuntu-latest
    steps:
      - uses: actions/checkout@v4

      - name: Setup Java 17 (Eclipse Temurin)
        uses: actions/setup-java@v4
        with:
          java-version: '17'
          distribution: 'temurin'
          cache: 'maven'

      - name: Run Tests & Build JAR
        run: cd backend && mvn clean package

      - name: Upload Backend JAR
        uses: actions/upload-artifact@v4
        with:
          name: backend-jar
          path: backend/target/*.jar
  \end{lstlisting}
\end{frame}

\begin{frame}[fragile]{Paso 3: Implementar Job 2 (\texttt{test-frontend})}
  \vspace{0.05cm}
  Responsable de validar tipado TypeScript y empaquetar la SPA en React:
  \vspace{0.1cm}
  \begin{lstlisting}[style=yamltiny]
  test-frontend:
    name: Frontend Lint & Build
    runs-on: ubuntu-latest
    steps:
      - uses: actions/checkout@v4

      - name: Setup Node.js 20
        uses: actions/setup-node@v4
        with:
          node-version: 20
          cache: 'npm'
          cache-dependency-path: frontend/package-lock.json

      - name: Install & Run Lint
        run: |
          cd frontend
          npm install
          npm run lint

      - name: Build SPA Bundle
        run: cd frontend && npm run build
  \end{lstlisting}
\end{frame}

\begin{frame}[fragile]{Paso 4: Implementar Job 3 (\texttt{build-docker})}
  \vspace{0.05cm}
  Construye las imágenes Docker utilizando \textit{Docker Buildx}:
  \vspace{0.1cm}
  \begin{lstlisting}[style=yamltiny]
  build-docker:
    name: Build Multi-Stage Docker Images
    needs: [test-backend, test-frontend]
    runs-on: ubuntu-latest
    steps:
      - uses: actions/checkout@v4

      - name: Set up Docker Buildx
        uses: docker/setup-buildx-action@v3

      - name: Build Backend Docker Image
        uses: docker/build-push-action@v5
        with:
          context: ./backend
          push: false
          tags: product-api:latest

      - name: Build Frontend Docker Image
        uses: docker/build-push-action@v5
        with:
          context: ./frontend
          push: false
          tags: product-frontend:latest
  \end{lstlisting}
\end{frame}

\begin{frame}[fragile]{Paso 5: Implementar Job 4 (\texttt{smoke-test})}
  \vspace{-0.1cm}
  Despliega el stack multicontenedor y valida su salud con pruebas de humo:
  \vspace{0.05cm}
  \begin{lstlisting}[style=yamltiny]
  smoke-test:
    name: Compose Up & Smoke Testing
    needs: [build-docker]
    runs-on: ubuntu-latest
    steps:
      - uses: actions/checkout@v4
      - name: Start Stack with Docker Compose
        run: docker compose up -d --build
      - name: Wait for Backend Healthcheck
        run: sleep 20
      - name: Execute Smoke Test (cURL API)
        run: |
          STATUS=$(curl -s -o /dev/null -w "%{http_code}" http://localhost:8080/api/products)
          echo "HTTP Status: $STATUS"
          if [ "$STATUS" -ne 200 ]; then exit 1; fi
      - name: Stop Containers
        if: always()
        run: docker compose down
  \end{lstlisting}
\end{frame}

% ==============================================================================
% CONCLUSIÓN Y CIERRE
% ==============================================================================
\contentSlideQuoteStat
  {Conclusión del Aprendizaje}
  {100\%}
  {Despliegues Reproducibles y Automatizados}
  {"Un pipeline de CI/CD bien estructurado convierte el despliegue de software en una tarea rutinaria, predecible y segura, eliminando la incertidumbre de los despliegues manuales."}

\begin{frame}{Preguntas de Reflexión y Puesta en Marcha}
  \vspace{0.2cm}
  \begin{enumerate}
    \item \textbf{Fallos Rápidos:} ¿Por qué es crucial ejecutar las pruebas unitarias antes de comenzar la construcción de imágenes Docker?
    \item \textbf{Caché de Dependencias:} ¿Qué beneficio ofrece el uso de \texttt{cache: 'maven'} y \texttt{cache: 'npm'} en el tiempo de ejecución del runner?
    \item \textbf{Self-Hosted vs. Cloud:} ¿En qué escenario de IasLab convendría usar un runner en \texttt{grid102} en lugar de \texttt{ubuntu-latest}?
    \item \textbf{Pruebas de Humo:} ¿Por qué validar un código HTTP 200 con \texttt{curl} nos da mayor certeza que solo verificar que el contenedor esté corriendo?
  \end{enumerate}
  \vspace{0.3cm}
  \centering
  \textbf{\color{cbBlue}¡Éxitos en el desarrollo del Taller Evaluativo! (Tiempo: 2 Horas)}
\end{frame}

\end{document}
